\documentclass{article}
\usepackage[utf8]{inputenc}
\usepackage{amsmath, amssymb}
\usepackage{geometry}
\geometry{a4paper, margin=1in}

\title{TDA Wasserstein Classifier}
\author{Algebraic Topology Laboratory}
\date{}

\begin{document}
\maketitle

\section*{Problem Description}

In this optional laboratory, you will use the \texttt{gudhi} library to build a robust shape classifier based on Persistent Homology.

You are given three ``clean'' reference point clouds: a circle ($H_1$ Betti number $\beta_1=1$), a figure-eight ($\beta_1=2$), and a Gaussian cluster ($\beta_1=0$). You are also given a noisy, randomly rotated point cloud, and you must determine which of the three reference shapes it was sampled from.

Since the point clouds themselves cannot be easily compared (due to noise and rotation), you will compare their \textit{Persistence Diagrams}. You will use the \textbf{Wasserstein Distance} ($W_1$) to quantify the similarity between the diagrams. The noisy shape will be classified as the reference shape that minimizes the Wasserstein distance.

\section*{Implementation Tasks}

Open \texttt{implementation\_tasks.py} and implement the following functions:

\subsection*{1. \texttt{compute\_h1\_diagram(point\_cloud, max\_edge\_length)}}
Use \texttt{gudhi.RipsComplex} to build a Vietoris-Rips filtration of the \texttt{point\_cloud}. Create a simplex tree up to dimension 2, and compute its persistence. Filter the resulting diagram to return only the $H_1$ intervals as an $N \times 2$ array of \texttt{[birth, death]} values.

\subsection*{2. \texttt{classify\_shape(noisy\_diag, reference\_diags)}}
You are provided the $H_1$ persistence diagram of the noisy shape, and a dictionary of the reference diagrams: \texttt{\{"circle": diag\_c, "figure8": diag\_f, "cluster": diag\_n\}}. 

Iterate through the dictionary and compute the Wasserstein distance using \texttt{gudhi.wasserstein.wasserstein\_distance} (with \texttt{order=1} and \texttt{internal\_p=2}). Return the string key (e.g., \texttt{"circle"}) that corresponds to the minimal distance.

\section*{Testing Your Code}
Run \texttt{python lab\_dashboard.py} to open the interactive visualization. 
\begin{itemize}
    \item Use the \textbf{Shape Explorer} to visualize how noise deforms the $H_1$ persistence diagram. The UI will draw exact matching lines between the clean and noisy diagrams to visualize the Wasserstein distance computation.
    \item Click \textbf{Run ML Evaluation} to generate a dataset of 300 noisy, rotated point clouds. Your classifier will run on the entire dataset and output a Confusion Matrix demonstrating its accuracy!
\end{itemize}

\end{document}
